\makeatletter
\def\input@path{{./}{../}{../../}{../../../}{../../../../}{preamble/}{../preamble/}{../../preamble/}{../../../preamble/}{../../../../preamble/}}
\makeatother
\input{preamble_exams.tex}
\input{graphs/G11_Histograms.tex}
\title{Grade 11 -- Term 1 \\ One Dataset, Many Representations: C6 and C7}
\author{}
\date{}
\pagestyle{empty}
\definecolor{IntervalColor}{HTML}{175A91}
\definecolor{MidpointColor}{HTML}{8A4500}
\definecolor{FrequencyColor}{HTML}{006B58}
\definecolor{FormulaColor}{HTML}{663399}
\definecolor{ResultColor}{HTML}{145A32}
\newcommand{\iv}[1]{\textcolor{IntervalColor}{#1}}
\newcommand{\mv}[1]{\textcolor{MidpointColor}{#1}}
\newcommand{\fq}[1]{\textcolor{FrequencyColor}{#1}}
\newcommand{\fm}[1]{\textcolor{FormulaColor}{#1}}
\newcommand{\su}[1]{\textcolor{black!75}{#1}}
\newcommand{\resultbox}[1]{\par\noindent\fcolorbox{ResultColor}{ResultColor!7}{\parbox{\dimexpr\linewidth-2\fboxsep-2\fboxrule\relax}{#1}}\par}
\newcommand{\pageheading}[1]{\begin{center}{\large\bfseries #1}\end{center}\vspace{-4mm}}
\begin{document}
% Opening and solution geometry, typography and empty page style follow template_print.tex.
\newgeometry{left=0.4in,right=0.4in,top=0.35in,bottom=0.35in}
\vspace*{-15mm}
\pageheading{Grade 11 -- Term 1: One Dataset, Many Representations}
\textbf{One worked example: download times (C6 and C7).}
A school technician records the download times of 20 learning files, in seconds.
Only the histogram survives; individual times are unavailable. Convert this one
sample into tables and a representative list, then estimate its mean, sample
variance and standard deviation. A supplied theoretical model for the process has
\[
\mu=10\text{ s},\qquad \sigma^2=25\text{ s}^2,\qquad P(5\le X<15)=0.80.
\]
These are model benchmarks, not results calculated from the sample. The grouped
data and probability benchmark follow the download-time context in the reference
exam; the mean and variance benchmarks are supplied for this extended example.
\begin{center}\small\bfseries
Histogram $\longrightarrow$ Interval Table $\longrightarrow$ Midpoint Frequency Table\\
$\longrightarrow$ List of Observations $\longrightarrow$ Statistical Calculations
\end{center}
\textbf{Color key:} \iv{intervals}; \mv{midpoints}; \fq{frequencies and counts};
\fm{formulas}; \su{substitutions}; green boxes for final results.
Labels retain these distinctions when printed in grayscale.
\section*{1. Read the histogram}
\vspace{-4mm}
% Reuse existing explicit rectangles, boundary ticks and frequency grid.
\HistogramGraph{(0,5),(5,10),(10,15),(15,20)}{4,8,6,2}
{0.62\textwidth}{5cm}{IntervalColor!20}{5}{1}
\begin{center}\small
Horizontal axis: \iv{$x=$ download time (seconds)}.\quad
Vertical axis: \fq{frequency (number of files)}.
\end{center}
Each bar spans five seconds. Read its height from the frequency axis:
\fq{4, 8, 6, 2}, from left to right. Classes include their left endpoint and
exclude their right endpoint: 5 seconds belongs to $[5,10)$, not $[0,5)$.
\section*{2. Convert to an interval frequency table}
\vspace{-4mm}
Copy each pair of horizontal boundaries and its bar height. No approximation
is needed for this transfer.
\begin{center}
\begin{tabularx}{0.90\textwidth}{|C{1cm}|C{3cm}|C{2cm}|X|}\hline
Class $i$ & \iv{Interval (s)} & \fq{$f_i$} & Reading from the histogram\\\hline
1 & \iv{$[0,5)$} & \fq{4} & Boundaries 0 and 5; height 4.\\\hline
2 & \iv{$[5,10)$} & \fq{8} & Boundaries 5 and 10; height 8.\\\hline
3 & \iv{$[10,15)$} & \fq{6} & Boundaries 10 and 15; height 6.\\\hline
4 & \iv{$[15,20)$} & \fq{2} & Boundaries 15 and 20; height 2.\\\hline
\end{tabularx}
\end{center}
\[
\fm{n=\sum_{i=1}^{r}f_i},\qquad \su{r=4,\quad n=4+8+6+2=20}.
\]
There are $r=4$ classes and $n=20$ observations. Frequencies count files;
intervals describe possible times. Neither is a statistical average.
\newpage
\restoregeometry
\vspace*{-15mm}
\pageheading{Worked Solution -- From Intervals to Observations}
\section*{3. Construct the midpoint frequency table}
For lower boundary $L_i$ and upper boundary $U_i$, use the formula
\[
\fm{m_i=\frac{L_i+U_i}{2}}.
\]
The midpoint lies halfway between the boundaries. Keep the frequency unchanged.
\begin{center}
\begin{tabular}{|c|c|c|c|}\hline
$i$ & \iv{Interval (s)} & \mv{Midpoint calculation (s)} & \fq{$f_i$}\\\hline
1 & \iv{$[0,5)$} & \mv{$(0+5)/2=2.5$} & \fq{4}\\\hline
2 & \iv{$[5,10)$} & \mv{$(5+10)/2=7.5$} & \fq{8}\\\hline
3 & \iv{$[10,15)$} & \mv{$(10+15)/2=12.5$} & \fq{6}\\\hline
4 & \iv{$[15,20)$} & \mv{$(15+20)/2=17.5$} & \fq{2}\\\hline
\end{tabular}
\end{center}
The resulting discrete frequency table is
\[
\begin{array}{c|rrrr}
\mv{m_i\text{ (s)}}&\mv{2.5}&\mv{7.5}&\mv{12.5}&\mv{17.5}\\\hline
\fq{f_i}&\fq{4}&\fq{8}&\fq{6}&\fq{2}
\end{array}
\]
For example, the four original times in $[0,5)$ could be anywhere within that
interval. We represent each by 2.5 seconds. This does not establish that any
actual measurement was 2.5, or that the times were uniformly spread in the class.
\resultbox{\textbf{Approximation begins here.} Midpoint replacement preserves
all 20 counts but discards variation within intervals. Its statistics are exact
for the representative data and estimates for the original observations.}
\section*{4. Expand into an explicit list of observations}
Let $y_j$ denote representative observation $j$, for $j=1,\ldots,n$.
Write each midpoint $m_i$ exactly $f_i$ times:
\[
\begin{aligned}
(y_1,\ldots,y_{20})=(&\underbrace{\mv{2.5,\ 2.5,\ 2.5,\ 2.5}}_{\fq{4\text{ observations}}},\\
&\underbrace{\mv{7.5,\ 7.5,\ 7.5,\ 7.5,\ 7.5,\ 7.5,\ 7.5,\ 7.5}}_{\fq{8\text{ observations}}},\\
&\underbrace{\mv{12.5,\ 12.5,\ 12.5,\ 12.5,\ 12.5,\ 12.5}}_{\fq{6\text{ observations}}},\\
&\underbrace{\mv{17.5,\ 17.5}}_{\fq{2\text{ observations}}}).
\end{aligned}
\]
The length is $4+8+6+2=20$. This sorted representative list has exactly the
midpoint table's frequencies; it does not recover the original times or recording order.

\textbf{Notation.} Write $\bar y$, $s_y^2$, and $s_y$ for the representative
list's mean, sample variance, and sample standard deviation. Write $\bar x$,
$s_x^2$, and $s_x$ for the unknown statistics of the original observations.
Thus $\bar x\approx\bar y$, $s_x^2\approx s_y^2$, and $s_x\approx s_y$.
The index $i$ labels four classes; $j$ labels twenty representative observations.
\newpage
\vspace*{-15mm}
\pageheading{Worked Solution -- C7: Descriptive Statistics}
\begingroup
\setlength{\parskip}{3pt}
\setlength{\abovedisplayskip}{6pt}
\setlength{\belowdisplayskip}{6pt}
\section*{5. Sample mean}\vspace{-4mm}
First state the list and equivalent weighted formulas:
\[
\fm{\bar y=\frac{\sum_{j=1}^{n}y_j}{n}
=\frac{\sum_{i=1}^{r}f_i m_i}{n}}.
\]
The product $f_i m_i$ adds midpoint $m_i$ a total of $f_i$ times.
\begin{center}
\begin{tabular}{|c|c|c|}\hline
\mv{$m_i$ (s)} & \fq{$f_i$} & \su{$f_i m_i$ (s)}\\\hline
\mv{2.5}&\fq{4}&$4(2.5)=10$\\\hline
\mv{7.5}&\fq{8}&$8(7.5)=60$\\\hline
\mv{12.5}&\fq{6}&$6(12.5)=75$\\\hline
\mv{17.5}&\fq{2}&$2(17.5)=35$\\\hline
Total&\fq{20}&180\\\hline
\end{tabular}
\end{center}
Substitute the midpoints and counts; divide by $n=20$, not by four classes:
\[
\su{\bar y=\frac{4(2.5)+8(7.5)+6(12.5)+2(17.5)}{20}
=\frac{10+60+75+35}{20}=9\text{ s}}.
\]
Summing the expanded list also gives 180 seconds, so both representations give 9 seconds.
\section*{6. Sample variance}\vspace{-4mm}
Use the reference convention: the sample denominator is $n-1$.
\[
\fm{s_y^2=\frac{\sum_{j=1}^{n}(y_j-\bar y)^2}{n-1}
=\frac{\sum_{i=1}^{r}f_i(m_i-\bar y)^2}{n-1}}.
\]
Subtract the unrounded mean 9 from each midpoint, square the difference, then
multiply by the frequency to include every observation.
\begin{center}
\begin{tabular}{|c|c|c|c|c|}\hline
\mv{$m_i$}&\fq{$f_i$}&$m_i-9$&$(m_i-9)^2$&$f_i(m_i-9)^2$\\\hline
\mv{2.5}&\fq{4}&$-6.5$&42.25&$4(42.25)=169$\\\hline
\mv{7.5}&\fq{8}&$-1.5$&2.25&$8(2.25)=18$\\\hline
\mv{12.5}&\fq{6}&3.5&12.25&$6(12.25)=73.5$\\\hline
\mv{17.5}&\fq{2}&8.5&72.25&$2(72.25)=144.5$\\\hline
\end{tabular}
\end{center}
\[
\su{\sum f_i(m_i-9)^2=169+18+73.5+144.5=405\text{ s}^2},
\qquad\su{s_y^2=\frac{405}{20-1}=\frac{405}{19}\text{ s}^2}.
\]
The list repeats these same squared deviations with the same frequencies, giving
405 as well. The denominator is 19, not the number of classes minus one.
\section*{7. Sample standard deviation}\vspace{-4mm}
\[
\fm{s_y=\sqrt{s_y^2}},\qquad
\su{s_y=\sqrt{\frac{405}{19}}\text{ s}\approx4.617\text{ s}}.
\]
\resultbox{\textbf{Midpoint estimates for the original sample:}
$\bar x\approx9.000$ s; $s_x^2\approx21.316$ s$^2$; $s_x\approx4.617$ s.
Mean and standard deviation have time units; variance has squared time units.}
\endgroup
\newpage
\vspace*{-15mm}
\pageheading{Worked Solution -- Errors Relative to the Model}
\section*{8. Absolute and percentage errors for C7}
For estimate $\widehat\theta$ and nonzero theoretical benchmark $\theta$, first use
\[
\fm{E_\theta=|\widehat\theta-\theta|},\qquad
\fm{R_\theta=\frac{E_\theta}{|\theta|}\times100\%}.
\]
Absolute error is the size of the difference, in the statistic's units.
Percentage error compares that difference with the theoretical value, not the
sample estimate. All benchmarks here are nonzero. First find the theoretical
standard deviation:
\[
\fm{\sigma=\sqrt{\sigma^2}},\qquad\su{\sigma=\sqrt{25}=5\text{ s}}.
\]
\textbf{Mean:} substitute $\bar y=9$ for the estimate and $\mu=10$ for the benchmark.
\[
\su{E_\mu=|9-10|=|-1|=1\text{ s}},\qquad
\su{R_\mu=\frac{1}{10}\times100\%=10.00\%}.
\]
\textbf{Variance:} substitute the exact fraction $405/19$ and $\sigma^2=25$.
\[
\begin{aligned}
\su{E_{\sigma^2}}&\su{=\left|\frac{405}{19}-25\right|
=\left|\frac{405-475}{19}\right|=\frac{70}{19}\text{ s}^2
\approx3.684\text{ s}^2},\\
\su{R_{\sigma^2}}&\su{=\frac{70/19}{25}\times100\%
=\frac{280}{19}\%\approx14.74\%}.
\end{aligned}
\]
\textbf{Standard deviation:} substitute $\sqrt{405/19}$ and $\sigma=5$.
Because $405/19<25$, its square root is less than 5; hence
\[
\begin{aligned}
\su{E_\sigma}&\su{=\left|\sqrt{\frac{405}{19}}-5\right|
=\left(5-\sqrt{\frac{405}{19}}\right)\text{ s}\approx0.383\text{ s}},\\
\su{R_\sigma}&\su{=\frac{5-\sqrt{405/19}}{5}\times100\%
\approx7.66\%}.
\end{aligned}
\]
Keep the exact fraction and square root in error calculations. Round only the
final statistics and absolute errors to three decimal places and percentage
errors to two decimal places.
\begin{center}
\begin{tabular}{|l|c|c|c|c|}\hline
Statistic&Midpoint estimate&Theoretical value&Absolute error&Percentage error\\\hline
Mean (s)&9.000&10.000&1.000&$10.00\%$\\\hline
Variance (s$^2$)&21.316&25.000&3.684&$14.74\%$\\\hline
Standard deviation (s)&4.617&5.000&0.383&$7.66\%$\\\hline
\end{tabular}
\end{center}
\resultbox{All three midpoint statistics are below their model benchmarks.
These errors compare midpoint estimates with the model. They do not isolate the
unknown error caused by grouping: that requires the original measurements.
Sampling variation can also cause a sample to differ from its theoretical model.}
\newpage
\vspace*{-15mm}
\pageheading{Worked Solution -- C6 and Representation Checks}
\section*{9. Interval probability from the same dataset}
The event $A=\{5\le X<15\}$ contains two complete classes,
\iv{$[5,10)$} and \iv{$[10,15)$}. Let $k$ count files in the event.
\[
\fm{k=\sum_{i:\,[L_i,U_i)\subseteq A}f_i},\qquad
\fm{\widehat P(A)=\frac{k}{n}}.
\]
Substitute the heights of the included bars and the total count:
\[
\su{k=8+6=14},\qquad\su{n=4+8+6+2=20},\qquad
\su{\widehat P(5\le X<15)=\frac{14}{20}=0.70}.
\]
The interval table gives the same count. The list contains eight copies of 7.5
and six copies of 12.5 in the event: again 14 out of 20. Midpoint replacement
changes no memberships here because the event contains whole classes.
Thus 0.70 is the exact relative frequency of the original sample, while it is
an empirical estimate of the process probability.

\textbf{Compare with the theoretical probability $p=0.80$.} First state
\[
\fm{E_p=|\widehat P(A)-p|},\qquad
\fm{R_p=\frac{E_p}{|p|}\times100\%}.
\]
Then substitute the sample proportion and the supplied benchmark:
\[
\su{E_p=|0.70-0.80|=0.10},\qquad
\su{R_p=\frac{0.10}{0.80}\times100\%=12.50\%}.
\]
\resultbox{\textbf{C6 result:} empirical probability $0.70$ (70\%);
absolute error $0.10$, or 10 percentage points; relative percentage error
$12.50\%$. Percentage points and percentage error are different comparisons.}
\section*{10. Check what each representation preserves}
\begin{center}
\begin{tabularx}{\textwidth}{|C{3.7cm}|X|}\hline
Representation&Information and checks\\\hline
Histogram and interval table&Same boundaries and frequencies; $n=20$.
Individual times are unknown.\\\hline
Midpoint frequency table&Four representative values with frequencies $4,8,6,2$.
Count stays 20; original times are approximated.\\\hline
Expanded list&Exactly the same representative data as the midpoint table:
20 entries, sum 180 s, squared-deviation sum 405 s$^2$.\\\hline
Statistical calculations&List and weighted formulas agree:
$\bar y=9$, $s_y^2=405/19$, $s_y=\sqrt{405/19}$.\\\hline
Interval probability&$14/20$ is exact for this sample because the event uses
whole classes. A boundary cutting through a class would require more information
or an extra approximation.\\\hline
\end{tabularx}
\end{center}
\textbf{Interpretation.} The midpoint data average 9 seconds per download with
sample standard deviation about 4.617 seconds. These estimate the original
sample's typical time and spread; model benchmarks describe the process.
Changing representation does not create additional observations.
\end{document}

